\section{Methodology}

\subsection{Simulated process and study design}

A latent state $z_t\in\mathbb{R}^8$ follows a stable mean-adjusted VAR(1), and the four observed
variables are a nonlinear function of it:
\begin{equation}
z_t=\mu_t+\Phi(z_{t-1}-\mu_{t-1})+\eta_t,\quad \eta_t\sim N(0,0.2I_8);\qquad
x_t=W_2\tanh(z_t)+\varepsilon_t,\quad \varepsilon_t\sim N(0,0.05I_4).
\label{eq:process}
\end{equation}
$\Phi$ is symmetric tridiagonal (diagonal 0.55, adjacent off-diagonals 0.15; spectral radius
$\approx0.832$), $\tanh$ acts componentwise and $W_2$ is a fixed $4\times8$ mixing matrix
(Appendix~\ref{app:model}). Each trajectory starts at $z_0=0$ and discards 500 burn-in steps.
Because the map is nonlinear and $p=4<8$, no linear projection of $x_t$ recovers $z_t$; the latent
state is never supplied to any model. The fault is a sustained latent mean shift
$\mu_t=\delta v$ in the fixed direction $v=(1,1,0,\dots,0)\T/\sqrt2$, starting at the first
retained observation ($\tau=1$), so $\delta$ is the Euclidean size of the shift and every
out-of-control trajectory is post-change from its first observation. Fifteen magnitudes,
$\delta=0.2,0.4,\dots,3.0$, are studied; no other fault type is used.

Five datasets generated from independent random streams have separate roles
(Figure~\ref{fig:pipeline}): D1--D3 fit and calibrate the scheme in Phase~I, after which everything is
frozen and applied unchanged to D4 (independent in-control validation) and D5 (out-of-control
evaluation) in Phase~II.

\begin{figure}[!t]
\centering
\includegraphics[width=\linewidth]{methodology_pipeline_diagram.pdf}
\caption[Monitoring pipeline and data roles]{Hybrid and benchmark monitoring paths (top) and the
roles of the five datasets (bottom). Box shading marks the dataset on which a stage is fitted or
calibrated; every chart first reports at raw time $t=32$ and has $J=2969$ opportunities per
trajectory.}
\label{fig:pipeline}
\end{figure}

\subsection{Learned components}

\textbf{Windows.} D1 is split chronologically into training rows $[0,1600)$ and validation rows
$[1600,2000)$ before any window is formed. Each variable is standardised with the training-segment
mean and standard deviation, which are then frozen. Windows $X_t\in\mathbb{R}^{4\times32}$ of the
$L=32$ most recent standardised observations are formed with stride 1 inside a segment or
trajectory, so a 3000-observation trajectory yields $J=2969$ windows, the first ending at $t=32$.

\textbf{CAE.} The encoder applies two one-dimensional convolutions ($4\to8\to16$ channels,
kernel 3, stride 2, ReLU), which halve the window length twice ($32\to16\to8$), and a linear
bottleneck producing $h_t\in\mathbb{R}^8$; the decoder mirrors this with a linear layer and two
transposed convolutions (3180 parameters; Appendix~\ref{app:model}). The network was trained on
the 1569 D1 training windows to minimise the mean squared reconstruction error (Adam, learning rate
$10^{-3}$, batch size 64, at most 100 epochs, early-stopping patience 10), and the checkpoint with
the lowest validation loss (epoch 33) was frozen. The reconstruction component is the summed
squared prediction error of the standardised window and its floored logarithm,
\begin{equation}
\SPE_t=\sum_{c=1}^{4}\sum_{k=1}^{32}\big(X_{t,ck}-\hat X_{t,ck}\big)^2,\qquad
\ell_t=\log\max(\SPE_t,10^{-12}).
\label{eq:spe}
\end{equation}

\textbf{OC-SVM.} The 1969 D2 windows are passed through the frozen standardisation and encoder, and
an RBF OC-SVM is fitted to the resulting features directly, without further scaling:
\begin{equation}
k(h,h')=\exp\!\big(-\gamma\lVert h-h'\rVert^2\big),\qquad \nu=0.03,\qquad
\gamma=\Big[\operatorname*{median}_{i<j}\lVert h_i-h_j\rVert^2\Big]^{-1}\approx0.0305,
\label{eq:ocsvm}
\end{equation}
with the median taken over all $\binom{1969}{2}$ pairs. The parameter $\nu$ bounds the fraction of
training windows outside the fitted boundary from above and the fraction of support vectors from
below \citep{scholkopf_estimating_2001}; the fit has 78 support vectors. The monitored score is the
negated decision function, $s_t=\rho-\sum_j\alpha_j k(h_{(j)},h_t)$, so that a larger $s_t$ means a
more abnormal window.

\subsection{Monitoring charts and calibration}

The hybrid monitoring vector is $y_t=(s_t,\ell_t)\T$. For a chart input $u_j$ with in-control mean $\hat\mu$ and covariance $\hat\Sigma$, the MEWMA state
and statistic are \citep{lowry_multivariate_1992}
\begin{equation}
Z_j=\lambda u_j+(1-\lambda)Z_{j-1},\;\; Z_0=\hat\mu,\qquad
T_j^2=(Z_j-\hat\mu)\T(c_j\hat\Sigma)^{-1}(Z_j-\hat\mu),\;\;
c_j=\tfrac{\lambda}{2-\lambda}\big[1-(1-\lambda)^{2j}\big],
\label{eq:mewma}
\end{equation}
where $j$ indexes monitoring opportunities ($j=1$ at $t=32$) and the exact finite-start factor
$c_j$ is always used. Six configurations of four charts are evaluated (Table~\ref{tab:ic}). The
hybrid chart $H$ monitors $y_t$ at the primary $\lambda=0.05$, with $\lambda=0.10$ and $0.20$ as a
sensitivity analysis; the ablations $S$ and $R$ monitor $s_t$ and $\ell_t$ alone, and the benchmark
$M$ is a conventional MEWMA on the raw $x_t$, each with $\lambda=0.05$. $\hat\mu$ and
$\hat\Sigma$ are pooled over all $300\times2969$ D3 monitoring vectors ($S$ and $R$ use the matching
entries) and, for $M$, over the aligned raw D3 observations at $t=32,\dots,3000$. Because the hybrid
charts cannot decide before the first window is complete, $M$'s state is updated silently with
$x_1,\dots,x_{31}$ and it first reports at $t=32$ with factor $c_{32}$, so all six charts share the
same 2969 opportunities.

Consecutive windows share 31 of their 32 observations, so the monitored sequences are strongly
autocorrelated (Section~\ref{sec:components}) and $c_j\hat\Sigma$ is only a reference scaling, not
the covariance of $Z_j$; without independence, limits derived for independent data lose their
nominal ARL properties \citep{lowry_multivariate_1992}. Each limit is therefore calibrated
empirically on D3. For a candidate $h$, $\widehat{\ARL}_0(h)$ is the mean over the 300 D3
trajectories of the first opportunity with $T_j^2>h$, truncated at $J=2969$. As this is a
non-decreasing step function of $h$, the calibrated limit
\begin{equation}
h^*=\min\big\{h\in V:\widehat{\ARL}_0(h)\ge370\big\},
\label{eq:limit}
\end{equation}
where $V$ is the set of values attained by the running maxima of $T_j^2$, is found exactly by
bisection. The target of 370 is the $\ARL_0$ of a conventional three-sigma Shewhart chart
\citep{ahmed_statistical_2017}.

\subsection{Evaluation design}

A run length is the first opportunity $j$ with $T_j^2>h^*$; if none occurs by $j=2969$, the run
length is recorded as 2969 and flagged as censored by a separate indicator, so each mean is a
restricted ARL. D4 yields each chart's independent in-control run lengths, kept distinct from its D3
calibration value, and D5 its out-of-control run lengths at the 15 shifts. Because $\tau=1$, the first window is already post-change and a signal at
opportunity $j$ corresponds to $j+31$ post-change observations: $\ARL_1$ is a detection delay
measured from the first permitted decision, the setting in which \citet{lowry_multivariate_1992}
also compared multivariate charts. $H(0.05)$ and $M(0.05)$ monitor the same D5 trajectories, so the
primary comparison is paired,
\begin{equation}
\bar d(\delta)=\frac{1}{500}\sum_{i=1}^{500}\big[\mathrm{RL}_{H,i}(\delta)-\mathrm{RL}_{M,i}(\delta)\big],
\label{eq:paired}
\end{equation}
with standard error $s_d/\sqrt{500}$; a positive $\bar d$ means that the benchmark signalled first.
The design was fixed in advance: no model, moment, smoothing constant or limit was changed after D4
or D5 was examined.
